\documentclass[11pt]{article}

\usepackage[a4paper,margin=2.4cm]{geometry}
\usepackage{amsmath,amssymb,amsthm,mathtools}
\usepackage{array}
\usepackage{booktabs}
\usepackage{longtable}
\usepackage{enumitem}
\usepackage[colorlinks=true,linkcolor=blue,citecolor=blue,urlcolor=blue]{hyperref}
\usepackage{xcolor}
\usepackage[export]{adjustbox}
\usepackage{microtype}
\usepackage{seqsplit}
\usepackage{multirow}
\emergencystretch=2.5em

\title{The Toolbox: Category Rendering, Click-to-Add, and Drag-and-Drop\\
\large A Custom HTML Toolbox Over Blockly, and How a Block Enters the Workspace}
\author{Research note --- Block-MiniJava}
\date{2026-07-20}

\begin{document}
\maketitle

\begin{abstract}
Block-MiniJava does not use Blockly's built-in flyout-based toolbox. The
category list, search box, and every draggable block button in
\texttt{\#toolbox-column} are hand-rolled HTML
(\texttt{\seqsplit{src/core/ui/app.ts}}), driven by a small block/category registry
(\texttt{\seqsplit{src/core/blocks/minijavaBlocks.ts}}) and native HTML5
drag-and-drop. A block can enter the workspace two ways --- a click, or a
drag --- and both funnel through the same block-creation function. This
note documents the registry, the rendering, both entry paths (including
the exact client-to-workspace coordinate conversion drag-and-drop needs),
and the "required blocks" mechanism that keeps exactly one \texttt{Goal}
and one \texttt{MainClass} in the workspace at all times, self-healing
after any edit, and correspondingly \emph{absent} from the toolbox itself.
\end{abstract}

\tableofcontents

% =============================================================================
\section{Why a custom toolbox}

Blockly ships its own toolbox/flyout system (an XML or JSON toolbox
definition, rendered by Blockly itself as a sidebar flyout you drag out
of). Block-MiniJava does not use it. Instead, \texttt{\#toolbox-column}
(\texttt{\seqsplit{index.html:106--163}}) is ordinary HTML that \texttt{app.ts}
populates and wires by hand: a search input, then one collapsible
\texttt{<section class="toolbox-category">} per category, each holding a
\texttt{<button class="toolbox-block-button">} per block. This buys three
things a stock Blockly flyout does not give for free in this project's
IDE-shell design: the toolbox can live inside the same activity-bar/sidebar
system as the Run and Settings views (\texttt{\seqsplit{ui-design-and-layout.tex}}
\S2--\S3), the search box can filter by block \emph{label} text rather than
only by category, and every block button can be a real HTML
\texttt{draggable} element participating in the browser's native
drag-and-drop rather than Blockly's own (touch-oriented) flyout drag
gesture.

% =============================================================================
\section{The block and category registry}

\texttt{\seqsplit{src/core/blocks/minijavaBlocks.ts}} is the single source of truth
for which block types exist and how the toolbox groups them --- the
\texttt{renderToolbox} function (\S\ref{sec:render}) never hard-codes a
block type or category; it only iterates this data.

\begin{center}
\begin{longtable}{@{}p{4.0cm}p{10.2cm}@{}}
\toprule
\textbf{Symbol} & \textbf{Role} \\
\midrule
\endhead
\texttt{ToolboxBlock} &
  \texttt{\seqsplit{minijavaBlocks.ts:4--8}} --- \{type, label, icon\}: a block type
  name, its toolbox display label, and a short glyph/text icon. \\
\texttt{ToolboxCategory} &
  \texttt{\seqsplit{minijavaBlocks.ts:10--15}} --- \{id, label, icon, blocks:
  ToolboxBlock[]\}. \\
\texttt{\seqsplit{MINI\_JAVA\_REQUIRED\_BLOCK\_TYPES}} &
  \texttt{\seqsplit{minijavaBlocks.ts:17}} --- the two-element tuple
  \texttt{['mj\_goal', 'mj\_main\_class']}, destructured in \texttt{app.ts:62}
  into \texttt{\seqsplit{GOAL\_BLOCK\_TYPE}}/\texttt{\seqsplit{MAIN\_BLOCK\_TYPE}}. \\
\texttt{\seqsplit{MINI\_JAVA\_BLOCK\_TYPES}} &
  \texttt{\seqsplit{minijavaBlocks.ts:19--59}} --- the full list of 38 block type
  names (grammar non-terminals, statements, expressions, literals); the
  type \texttt{\seqsplit{MiniJavaBlockType}} (l.61) is its element union. \\
\texttt{\seqsplit{MINI\_JAVA\_CATEGORIES}} &
  \texttt{\seqsplit{minijavaBlocks.ts:67--144}} --- six \texttt{ToolboxCategory}
  entries, in display order: \textbf{Program}, \textbf{Declarations},
  \textbf{Types}, \textbf{Statements}, \textbf{Expressions},
  \textbf{Values} (\S\ref{sec:categories}). \\
\texttt{\seqsplit{defineMiniJavaBlocks()}} &
  \texttt{\seqsplit{minijavaBlocks.ts:155--}}\linebreak(end) --- registers every block's
  JSON definition with Blockly (\texttt{Blockly.common.}\linebreak\texttt{\seqsplit{defineBlocksWithJsonArray}}),
  including each input's grammar-aware connection \texttt{check} (\S\ref{sec:connections}) ---
  called once from \texttt{app.ts} before the workspace is injected. \\
\bottomrule
\end{longtable}
\end{center}

\subsection{The six categories}
\label{sec:categories}

\begin{center}
\begin{longtable}{@{}p{2.6cm}p{2.2cm}p{9.4cm}@{}}
\toprule
\textbf{Category} & \textbf{id} & \textbf{Blocks} \\
\midrule
\endhead
Program & \texttt{program} & \texttt{\seqsplit{mj\_class\_declaration}}. \\
Declarations & \texttt{declarations} & \texttt{\seqsplit{mj\_var\_declaration}},
  \texttt{\seqsplit{mj\_method\_declaration}}, \texttt{\seqsplit{mj\_formal\_parameter}}. \\
Types & \texttt{types} & \texttt{\seqsplit{mj\_type\_int\_array}},
  \texttt{\seqsplit{mj\_type\_boolean}}, \texttt{mj\_type\_int}, \texttt{mj\_type\_string},
  \texttt{\seqsplit{mj\_type\_identifier}}. \\
Statements & \texttt{statements} & \texttt{\seqsplit{mj\_statement\_block}},
  \texttt{\seqsplit{mj\_statement\_if}}, \texttt{\seqsplit{mj\_statement\_while}},
  \texttt{\seqsplit{mj\_statement\_print}}, \texttt{\seqsplit{mj\_statement\_assign}},
  \texttt{\seqsplit{mj\_statement\_array\_assign}}. \\
Expressions & \texttt{expressions} & \texttt{mj\_expr\_arith},
  \texttt{\seqsplit{mj\_expr\_compare}}, \texttt{mj\_expr\_logic},
  \texttt{\seqsplit{mj\_expr\_array\_lookup}}, \texttt{\seqsplit{mj\_expr\_array\_length}},
  \texttt{\seqsplit{mj\_expr\_char\_at}}, \texttt{mj\_expr\_concat},
  \texttt{\seqsplit{mj\_expr\_str\_length}}, \texttt{\seqsplit{mj\_expr\_method\_call}},
  \texttt{\seqsplit{mj\_argument\_item}}, \texttt{mj\_expr\_not},
  \texttt{mj\_expr\_parens}. \\
Values & \texttt{values} & \texttt{\seqsplit{mj\_expr\_integer}},
  \texttt{mj\_expr\_string}, \texttt{\seqsplit{mj\_expr\_boolean}},
  \texttt{\seqsplit{mj\_expr\_identifier}}, \texttt{mj\_expr\_this},
  \texttt{\seqsplit{mj\_expr\_new\_int\_array}}, \texttt{\seqsplit{mj\_expr\_new\_object}}. \\
\bottomrule
\end{longtable}
\end{center}

Two block types from \texttt{\seqsplit{MINI\_JAVA\_BLOCK\_TYPES}} are conspicuously
absent from every category: \texttt{mj\_goal} and \texttt{mj\_main\_class}
--- see \S\ref{sec:required} for why. \texttt{\seqsplit{mj\_value\_object}},
\texttt{mj\_value\_null}, and \texttt{\seqsplit{mj\_viz\_description}} are also absent:
the first two are runtime-only value blocks the substitution stepper
renders into its own private workspace (\texttt{\seqsplit{typing-and-evaluation-
rules.tex}} \S4.3, \texttt{\seqsplit{bottom-panel-sync.tex}} \S2.4), never
user-authored; the third is a visualization annotation block, also never
user-placed directly.

% =============================================================================
\section{Rendering the toolbox}
\label{sec:render}

\begin{center}
\begin{longtable}{@{}p{3.6cm}p{10.6cm}@{}}
\toprule
\textbf{Function} & \textbf{Role} \\
\midrule
\endhead
\texttt{\seqsplit{renderToolbox(query)}} &
  \texttt{app.ts:830--898} --- rebuilds \texttt{\seqsplit{\#toolbox-content}} from
  scratch on every call (no incremental diffing --- the whole list is
  cheap enough to throw away and rebuild). With no query, every category
  and block renders. With a query, a category whose own label/id matches
  shows \emph{all} its blocks; otherwise each block is individually
  filtered by whether its label or type contains the (lower-cased) query
  (l.837--843) --- so searching \texttt{"if"} shows the \texttt{if / else}
  statement even though its category label ("Statements") does not
  itself match. A category that ends up with zero visible blocks is
  omitted entirely (l.844); if \emph{no} block anywhere matches, a
  \texttt{\seqsplit{.toolbox-search-empty}} message is shown instead
  (l.891--897). \\
\texttt{\seqsplit{addBlockToWorkspace(type)}} &
  \texttt{app.ts:900--902} --- the click handler installed on every
  toolbox button (l.872); see \S\ref{sec:click}. \\
\bottomrule
\end{longtable}
\end{center}

Each category is a collapsible \texttt{<section>}: its
\texttt{<button class="toolbox-category-header">} toggles a
\texttt{collapsed} class on click (\texttt{app.ts:855--858}), independent
per category, with no persisted collapse state --- collapsing is a
same-session convenience, not a saved preference.

% =============================================================================
\section{Entry path 1: click to add}
\label{sec:click}

\begin{center}
\begin{longtable}{@{}p{4.2cm}p{10.0cm}@{}}
\toprule
\textbf{Function} & \textbf{Role} \\
\midrule
\endhead
\texttt{\seqsplit{addBlockToWorkspace(type)}} &
  \texttt{app.ts:900--902} --- \texttt{\seqsplit{createBlockInWorkspace(type,}}\linebreak\texttt{\seqsplit{activeWorkspaceInsertionPoint())}}. \\
\texttt{\seqsplit{activeWorkspaceInsertionPoint()}} &
  \texttt{app.ts:178--185} --- picks a point inset from the visible
  workspace viewport's top-left corner (18\% of the viewport's
  width/height, clamped to $[72,\allowbreak 140]$px horizontally and
  $[72,120]$px vertically, l.180--181), then adds a small
  \textbf{cascading offset} that increments by 24px (mod 120) on every
  call (\texttt{clickAddOffset}, l.50, l.182--183) --- so repeatedly
  clicking the same toolbox button drops a visible \emph{staircase} of
  blocks rather than stacking them exactly on top of each other. \\
\texttt{\seqsplit{createBlockInWorkspace(type, position)}} &
  \texttt{app.ts:187--205} --- if \texttt{type} is the Goal or MainClass
  type, redirects to \texttt{\seqsplit{ensureRequiredBlocks}} instead of creating a
  duplicate (a defensive branch: neither type is ever reachable through
  the toolbox UI at all, \S\ref{sec:required}, so this only guards a
  hypothetical direct call). Otherwise: opens a
  \texttt{Blockly.Events} group (l.193, so the create + move the caller
  performs together undo as one user action), calls
  \texttt{\seqsplit{workspace.newBlock(type)}}, \texttt{initSvg()}, \texttt{render()},
  \texttt{moveBy(position.x, position.y)}, and \texttt{select()}, then
  calls \texttt{updateCode()} and \texttt{saveAutosave()}. \\
\bottomrule
\end{longtable}
\end{center}

\texttt{\seqsplit{clientPointToWorkspacePoint}} (\S\ref{sec:coords}) is \emph{not}
used on this path --- a plain click has no client coordinate to convert,
only "somewhere sensible near the top-left of what's currently visible."

% =============================================================================
\section{Entry path 2: drag and drop}
\label{sec:dnd}

Every toolbox block button is a native HTML5 draggable element
(\texttt{button.draggable = true}, \texttt{app.ts:869}); dropping it onto
the Blockly canvas creates a block at the exact drop point. This is
ordinary browser drag-and-drop (the \texttt{DataTransfer} API), not any
Blockly-provided drag gesture.

\subsection{Starting the drag: the toolbox button}

\begin{center}
\begin{longtable}{@{}p{4.0cm}p{10.2cm}@{}}
\toprule
\textbf{Event/Code} & \textbf{Role} \\
\midrule
\endhead
\texttt{\seqsplit{TOOLBOX\_BLOCK\_MIME}} &
  \texttt{app.ts:61} --- the custom MIME type
  \texttt{\seqsplit{"application/x-block-minijava-block"}} used to carry the block
  \texttt{type} string through the drag. \\
\texttt{dragstart} &
  \texttt{app.ts:873--878} --- adds the \texttt{is-dragging} CSS class for
  visual feedback; calls \texttt{\seqsplit{event.dataTransfer.setData}} \textbf{twice}
  --- once under \texttt{\seqsplit{TOOLBOX\_BLOCK\_MIME}} (the precise channel) and
  once under the standard \texttt{'text/plain'} (a fallback, in case
  something in the drop path only reads the generic type); sets
  \texttt{effectAllowed = 'copy'} (dragging a toolbox block never removes
  it from the toolbox --- it is always a copy/instantiate operation, never
  a move). \\
\texttt{dragend} &
  \texttt{app.ts:879--882} --- removes \texttt{is-dragging} and clears the
  \texttt{\seqsplit{workspace-drop-target}} highlight regardless of whether the drop
  actually succeeded (a drag cancelled outside any valid target still
  fires \texttt{dragend}). \\
\bottomrule
\end{longtable}
\end{center}

\subsection{Receiving the drag: the workspace area}

\begin{center}
\begin{longtable}{@{}p{4.0cm}p{10.2cm}@{}}
\toprule
\textbf{Function} & \textbf{Role} \\
\midrule
\endhead
\texttt{\seqsplit{hasToolboxBlockDrag(event)}} &
  \texttt{app.ts:908--910} --- true iff the in-flight drag's
  \texttt{\seqsplit{dataTransfer.types}} includes \texttt{\seqsplit{TOOLBOX\_BLOCK\_MIME}} ---
  the guard every handler below uses first, so dragging some unrelated
  thing (a browser tab, a file, page text) over the canvas is silently
  ignored rather than mishandled. \\
\texttt{\seqsplit{initToolboxDragAndDrop()}} &
  \texttt{app.ts:912--941} --- installs four listeners on
  \texttt{\#blockly-area}: \\
  & \quad\texttt{dragenter} (l.915--919) --- \texttt{preventDefault()}
  (required for the element to become a valid drop target at all under
  the HTML5 DnD spec) and adds \texttt{\seqsplit{workspace-drop-target}} for a
  visible drop-zone highlight. \\
  & \quad\texttt{dragover} (l.921--926) --- also \texttt{preventDefault()}
  (fired continuously while hovering; the spec requires this every time,
  not just on \texttt{dragenter}), and sets \texttt{dropEffect = 'copy'}
  so the OS cursor shows a copy (not move) affordance. \\
  & \quad\texttt{dragleave} (l.928--931) --- removes the highlight only
  when the pointer has left \texttt{\#blockly-area} entirely (checked via
  \texttt{\seqsplit{event.relatedTarget}}, so moving between child elements
  \emph{inside} the drop zone does not flicker the highlight off and back
  on). \\
  & \quad\texttt{drop} (l.933--940) --- \texttt{preventDefault()} (stops
  the browser's default handling, e.g.\ navigating to dropped text as a
  URL), clears the highlight, reads the block type back out of
  \texttt{dataTransfer} (custom MIME first, \texttt{'text/plain'} as
  fallback), and calls \texttt{\seqsplit{addBlockToWorkspaceAtClientPoint(type,}}\linebreak\texttt{event.clientX, event.clientY)}. \\
\bottomrule
\end{longtable}
\end{center}

\subsection{Placing the block: client-to-workspace coordinates}
\label{sec:coords}

A \texttt{drop} event's \texttt{clientX}/\texttt{clientY} are
\textbf{screen-space pixel} coordinates; a Blockly block needs a position
in \textbf{workspace-space} coordinates (which account for the canvas's
current pan offset and zoom level). This conversion is the one genuinely
non-trivial piece of the whole toolbox implementation.

\begin{center}
\begin{longtable}{@{}p{4.0cm}p{10.2cm}@{}}
\toprule
\textbf{Function} & \textbf{Role} \\
\midrule
\endhead
\texttt{\seqsplit{blocklyViewportRect()}} &
  \texttt{app.ts:159--161} --- \texttt{\#blockly-div}'s
  \texttt{\seqsplit{getBoundingClientRect()}}: the on-screen rectangle the canvas
  currently occupies. \\
\texttt{\seqsplit{clientPointToWorkspacePoint(clientX, clientY)}} &
  \texttt{app.ts:163--176} --- the conversion, in four steps:
  (1) read \texttt{\seqsplit{workspace.getMetrics()}} for the current
  \texttt{viewLeft}/\texttt{viewTop} (how far the canvas is scrolled/panned)
  and \texttt{currentScale()} for the zoom factor; (2) clamp the client
  point into the viewport rectangle first (\texttt{Math.max(0, Math.min(...))},
  l.170--171) so a drop event that lands slightly outside the canvas
  element (a few pixels of border/padding) still resolves to a point
  \emph{inside} it, rather than a nonsensical negative or out-of-range
  workspace coordinate; (3)-(4) convert the clamped local pixel offset to
  workspace units: $x_{\mathrm{ws}} = \mathrm{viewLeft} +
  \mathrm{localX}/\mathrm{scale}$, and symmetrically for $y$
  (l.172--175) --- dividing by scale is what makes a drop land on the
  \emph{same visual spot} regardless of whether the canvas is zoomed in
  or out. \\
\texttt{\seqsplit{addBlockToWorkspaceAtClientPoint(type, clientX, clientY)}} &
  \texttt{app.ts:904--906} --- \texttt{\seqsplit{createBlockInWorkspace(type,}}\linebreak\texttt{clientPointToWorkspacePoint(clientX, clientY))};
  the drag-and-drop counterpart of
  \texttt{\seqsplit{addBlockToWorkspace}} (\S\ref{sec:click}), sharing the same
  \texttt{\seqsplit{createBlockInWorkspace}} block-creation logic --- the only
  difference between the click path and the drag path is \emph{which}
  position function computes where the new block lands. \\
\bottomrule
\end{longtable}
\end{center}

\texttt{currentScale()} (\texttt{app.ts:150--156}) reads Blockly's own
\texttt{getScale()} when available, falling back to a raw \texttt{.scale}
property read for older/differently-typed Blockly builds --- defensive
compatibility code, not part of the coordinate math itself.

% =============================================================================
\section{Grammar-aware connections: what happens after the block lands}
\label{sec:connections}

Neither the click nor the drag path checks \emph{where} the new block is
allowed to connect --- a freshly created block always lands as an
unconnected top-level block at the computed position. Whether it can
subsequently be dragged \emph{into} a connection with an existing block is
governed entirely by Blockly's own connection-compatibility check against
the \texttt{check} metadata every block input declares in
\texttt{\seqsplit{defineMiniJavaBlocks()}} (\texttt{\seqsplit{minijavaBlocks.ts}}, 40
\texttt{check:} sites plus typed \texttt{output}/\texttt{\seqsplit{previousStatement}}/
\texttt{nextStatement} declarations) --- named after MiniJava grammar
non-terminals (\texttt{MainClass}, \texttt{ClassDeclaration},
\texttt{VarDeclaration}, \texttt{\seqsplit{MethodDeclaration}},
\texttt{FormalParameter}, \texttt{Type}, \texttt{Statement},
\texttt{Expression}, \texttt{Identifier}), exactly as the README's Grammar
section and the project's custom \texttt{BMJ-Thrasos} renderer
(\texttt{core/renderer/}) describe. This note does not re-derive that
system --- it is the same static, structural check documented in the
README's "Grammar-aware renderer" section --- but it is worth stating
precisely how the toolbox interacts with it: \textbf{the toolbox creates
blocks; the connection-check system, entirely separate and already built
into Blockly plus this project's \texttt{check} metadata, decides where
they may snap.} No toolbox code inspects or enforces grammar.

% =============================================================================
\section{Required blocks: Goal and MainClass}
\label{sec:required}

Every MiniJava program has exactly one \texttt{Goal} (the program root)
wrapping exactly one \texttt{MainClass}. Rather than ask the user to place
these two blocks correctly (and keep them that way), Block-MiniJava
enforces their existence and wiring automatically, and --- as already
noted in \S\ref{sec:categories} --- never offers them in the toolbox at
all, so there is no user action that could create a \emph{second} Goal or
a disconnected MainClass in the first place; the enforcement machinery
below is a safety net against indirect ways a workspace could still end up
without them (loading an old/hand-edited save file, an incomplete
text-editor import, an aborted paste).

\begin{center}
\begin{longtable}{@{}p{4.0cm}p{10.2cm}@{}}
\toprule
\textbf{Function} & \textbf{Role} \\
\midrule
\endhead
\texttt{\seqsplit{lockRequiredBlock(block)}} &
  \texttt{app.ts:255--259} --- \texttt{\seqsplit{setDeletable(false)}},
  \texttt{\seqsplit{setMovable(false)}}, \texttt{\seqsplit{setCollapsed(false)}}: once a Goal or
  MainClass block exists, the user can edit its contents (attach classes,
  fill in the class name) but can neither delete it nor drag it away from
  its required position. \\
\texttt{\seqsplit{createRequiredBlock(ws, type, x, y)}} &
  \texttt{app.ts:261--267} --- the same
  \texttt{newBlock}/\texttt{initSvg}/\texttt{render}/\texttt{moveBy}
  sequence as \texttt{\seqsplit{createBlockInWorkspace}} (\S\ref{sec:click}), factored
  out since required blocks are created from a different call site with
  no click/drag position to reuse. \\
\texttt{\seqsplit{removeDuplicateRequiredBlock(block)}} &
  \texttt{app.ts:269--272} --- \texttt{\seqsplit{setDeletable(true)}} then
  \texttt{dispose(false)}: a \emph{second} Goal or MainClass block (e.g.\
  from a merge of two saved files) is un-locked just long enough to be
  disposed of, keeping exactly one of each. \\
\texttt{\seqsplit{ensureRequiredBlocks(ws)}} &
  \texttt{app.ts:274--327} --- the sweep, guarded against re-entrancy
  (\texttt{\seqsplit{enforcingRequiredBlocks}}, l.51, l.275--276, since disposing a
  block or moving a connection inside this function itself fires more
  workspace-change events that would otherwise re-trigger it): finds (or
  creates) the one \texttt{Goal}; disposes every extra \texttt{Goal};
  finds (or creates) the \texttt{MainClass} \emph{connected to that Goal}
  in preference to any other stray \texttt{MainClass} on the workspace
  (l.297--298); disposes every other \texttt{MainClass}; if the Goal and
  MainClass are not already connected to each other, forcibly disconnects
  whatever \emph{is} on each connection first and connects them
  (l.312--319); finally locks both via \texttt{\seqsplit{lockRequiredBlock}}. Returns
  whether anything actually changed. \\
\texttt{\seqsplit{scheduleRequiredBlockEnforcement()}} &
  \texttt{app.ts:329--336} --- debounced (a zero-delay
  \texttt{setTimeout}, so it runs once after the current synchronous
  event-handling batch finishes rather than immediately inside a Blockly
  event callback) call to \texttt{\seqsplit{ensureRequiredBlocks}}; calls
  \texttt{updateCode()} only if something changed. Invoked from the
  workspace's change listener on every mutation event
  (\texttt{app.ts:1112,1121}) and from
  \texttt{\seqsplit{FINISHED\_LOADING}} (a full workspace load, e.g.\ Open or
  autosave restore, l.1111--1113). \\
\bottomrule
\end{longtable}
\end{center}

% =============================================================================
\section{Full inventory}

\begin{center}
\begin{longtable}{@{}p{4.4cm}p{5.0cm}p{5.4cm}@{}}
\toprule
\textbf{Path} & \textbf{Position function} & \textbf{Shared creation call} \\
\midrule
\endhead
Click a toolbox button &
  \texttt{\seqsplit{activeWorkspaceInsertionPoint()}} &
  \multirow{2}{5.4cm}{\texttt{\seqsplit{createBlockInWorkspace(type, position)}}} \\
Drag a toolbox button onto the canvas &
  \texttt{\seqsplit{clientPointToWorkspacePoint(}\linebreak\texttt{event.clientX, event.clientY)}} & \\
\bottomrule
\end{longtable}
\end{center}

% =============================================================================
\section{HTML wiring}

\begin{center}
\begin{longtable}{@{}p{3.8cm}p{10.4cm}@{}}
\toprule
\textbf{Element} & \textbf{Location} \\
\midrule
\endhead
\texttt{\#toolbox-column} &
  \texttt{\seqsplit{index.html:106--163}} --- the primary sidebar; the
  \texttt{blocks}/\texttt{search} view (\texttt{\seqsplit{index.html:114--120}})
  holds \texttt{\#toolbox-search} and \texttt{\seqsplit{\#toolbox-content}}. \\
\texttt{\#toolbox-search} &
  \texttt{index.html:116} --- wired at \texttt{app.ts:1219} to call
  \texttt{\seqsplit{renderToolbox(input.value)}} on every \texttt{input} event. \\
\texttt{\seqsplit{\#toolbox-content}} &
  \texttt{index.html:119} --- the render target of \texttt{renderToolbox}
  (\S\ref{sec:render}); empty until first populated at IDE init. \\
\texttt{\#blockly-area} &
  \texttt{\seqsplit{index.html:200--203}} --- the drop target for
  \texttt{\seqsplit{initToolboxDragAndDrop}} (\S\ref{sec:dnd}); wraps
  \texttt{\#blockly-div}, the actual Blockly-injected SVG canvas. \\
\bottomrule
\end{longtable}
\end{center}

% =============================================================================
\section*{Note on scope}
\addcontentsline{toc}{section}{Note on scope}
This note covers block \emph{creation} --- the toolbox UI, the click and
drag entry paths, and the required-block self-healing invariant. It does
not re-derive the grammar-aware connection-check system that governs where
a created block may subsequently connect (already documented in the
project's own README, "Grammar-aware renderer" section) or the static type
system that separately checks the resulting program
(\texttt{\seqsplit{typing-and-evaluation-rules.tex}}). The workbench shell the
toolbox sits inside (activity bar, sidebar resizing, responsive collapse)
is covered in \texttt{\seqsplit{ui-design-and-layout.tex}}.

\end{document}
